% >>> j2l compat: xcolor / named-colour compatibility
\PassOptionsToPackage{dvipsnames}{color}
\PassOptionsToPackage{dvipsnames,svgnames,x11names,table}{xcolor}
\RequirePackage{xcolor}
\makeatletter\@namedef{ver@color.sty}{2021/01/01}\makeatother

\makeatletter\@ifundefined{providecolor}{}{%
  \providecolor{Bittersweet}{RGB}{128,128,128}
  \providecolor{Dandelion}{RGB}{253,203,71}
  \providecolor{abstractcolor}{RGB}{128,128,128}
  \providecolor{rulecolor}{RGB}{128,128,128}
}\makeatother
% <<< j2l compat
\documentclass[10pt]{NSP1}
%%%%%%%%%%%%%%%%%%%%%%%%%%%%%%%%%%%%%%%%%%%%%%%%%%%%%%%%%%%%%%%%%%%%%%%%%%%%%%%%%%%%%%%%%%%%%%%%%%%%%%%%%%%%%%%%%%%%%%%%%%%%%%%%%%%%%%%%%%%%%%%%%%%%%%%%%%%%%%%%%%%%%%%%%%%%%%%%%%%%%%%%%%%%%%%%%%%%%%%%%%%%%%%%%%%%%%%%%%%%%%%%%%%%%%%%%%%%%%%%%%%%%%%%%%%%
\IfFileExists{url.sty}{\usepackage{url}}{}
\IfFileExists{floatflt.sty}{\usepackage{floatflt}}{}
\IfFileExists{helvet.sty}{\usepackage{helvet}}{}
\IfFileExists{times.sty}{\usepackage{times}}{}
\IfFileExists{psfig.sty}{\usepackage{psfig}}{}
\IfFileExists{graphics.sty}{\usepackage{graphics}}{}
\IfFileExists{mathptmx.sty}{\usepackage{mathptmx}}{}
\IfFileExists{amsmath.sty}{\usepackage{amsmath}}{}
\IfFileExists{amssymb.sty}{\usepackage{amssymb}}{}
\IfFileExists{bm.sty}{\usepackage{bm}}{}
\IfFileExists{float.sty}{\usepackage{float}}{}
\IfFileExists{caption.sty}{\usepackage[bf,hypcap]{caption}}{}
\IfFileExists{hyperref.sty}{\usepackage{hyperref}}{}
\IfFileExists{xcolor.sty}{\usepackage[dvipsnames,table,xcdraw]{xcolor}}{}
\IfFileExists{tabularx.sty}{\usepackage{tabularx}}{}
\IfFileExists{booktabs.sty}{\usepackage{booktabs}}{\providecommand\toprule{\hline}\providecommand\midrule{\hline}\providecommand\bottomrule{\hline}}
\IfFileExists{bm.sty}{\usepackage{bm}}{}
\IfFileExists{afterpage.sty}{\usepackage{afterpage}}{}
\IfFileExists{algorithm.sty}{\usepackage{algorithm}}{}
\IfFileExists{algpseudocode.sty}{\usepackage{algpseudocode}}{} % For pseudocode formatting
\IfFileExists{longtable.sty}{\usepackage{longtable}}{}


\setcounter{MaxMatrixCols}{10}
\addtolength{\oddsidemargin}{.5in}
\addtolength{\evensidemargin}{-.5in}

\tolerance=1
\emergencystretch=\maxdimen
\hyphenpenalty=10000
\hbadness=10000
\topmargin=0.00cm
\def\sm{\smallskip}
\def\no{\noindent}
\def\firstpage{1}
\setcounter{page}{\firstpage}
\def\thevol{}
\def\thenumber{}
\def\theyear{}

\usepackage{natbib}
\IfFileExists{graphicx.sty}{\usepackage{graphicx}}{}
\makeatletter
\def\biographyps#1#2{%
  \par\addvspace{21dd}\small\noindent
  \if!#1!\else
    \noindent\includegraphics[width=1in,height=1.25in,keepaspectratio]{#1}\par\smallskip
  \fi
  {\bfseries#2\unskip\ }\ignorespaces}
\def\endbiographyps{\par\addvspace{12pt}}
\makeatother
\IfFileExists{multirow.sty}{\usepackage{multirow}}{\providecommand\multirow[3]{#3}}
\IfFileExists{array.sty}{\usepackage{array}}{}
\IfFileExists{cuted.sty}{\usepackage{cuted}}{\newenvironment{strip}{\par\medskip\noindent\begin{minipage}{\linewidth}}{\end{minipage}\par\medskip}}
\makeatletter
\newcommand\jltx@wideopen{\par\addvspace{\medskipamount}\noindent}
\newcommand\jltx@wideclose{\par\addvspace{\medskipamount}}
\newenvironment{jltxspan}
  {\if@twocolumn\let\jltx@spannext\strip\else\let\jltx@spannext\jltx@wideopen\fi\jltx@spannext}
  {\if@twocolumn\let\jltx@spanend\endstrip\else\let\jltx@spanend\jltx@wideclose\fi\jltx@spanend}
\makeatother
\newcommand\JLTXaboveplacement{\par\addvspace{6pt plus 2pt minus 2pt}\noindent}
\newcommand\JLTXbelowplacement{\par\addvspace{6pt plus 2pt minus 2pt}}
\widowpenalty=10000
\clubpenalty=10000
\IfFileExists{newunicodechar.sty}{\usepackage{newunicodechar}}{\makeatletter\providecommand\newunicodechar[2]{%
  \begingroup\lccode`\~=`#1\lowercase{\endgroup\def~}{#2}%
  \catcode`#1=\active}\makeatother}
\newunicodechar{α}{\ensuremath{\alpha}}
\newunicodechar{÷}{\ensuremath{\div}}
\newunicodechar{≤}{\ensuremath{\leq}}

% >>> j2l compat: scalable Computer Modern (arbitrary font sizes)
\IfFileExists{type1cm.sty}{\usepackage{type1cm}}{}
\IfFileExists{fix-cm.sty}{\usepackage{fix-cm}}{}
% <<< j2l compat
\begin{document}
% >>> j2l compat: body-text size (+5%)
\makeatletter
\edef\JLTXbasesize{\f@size}
\@tempdimb=1.2\dimexpr\f@size pt\relax\edef\JLTXbaseskip{\strip@pt\@tempdimb}
\def\JLTXbodyscale{1.050}
\let\JLTX@normalsize\normalsize
\renewcommand\normalsize{%
  \JLTX@normalsize
  \@tempdima=\JLTXbodyscale\dimexpr\f@size pt\relax
  \@tempdimb=1.2\@tempdima
  \fontsize{\strip@pt\@tempdima}{\strip@pt\@tempdimb}\selectfont}%
\@ifundefined{@makecaption}{}{%
  \let\JLTX@makecaption\@makecaption
  \long\def\@makecaption#1#2{%
    \begingroup\fontsize{\JLTXbasesize}{\JLTXbaseskip}\selectfont
    \let\normalsize\JLTX@normalsize
    \JLTX@makecaption{#1}{#2}\endgroup}}%
\@ifundefined{thebibliography}{}{%
  \let\JLTX@thebibliography\thebibliography
  \def\thebibliography{%
    \fontsize{\JLTXbasesize}{\JLTXbaseskip}\selectfont\JLTX@thebibliography}}%
\makeatother
\normalsize
% <<< j2l compat





\titlefigurecaption{{\large \bf \rm Journal of Statistics Applications \& Probability }\\ {\it\small An International Journal}}

\title{Journal of Statistics Applications \& Probability}

\author{An International Journal \and Statistical Analysis of Social Responsibility Disclosure in Family Businesses \and Hanadi Abd-Almonem Ahmad AL-ahmar$^{1}$ \and Jafar Othman El- Sharif Abdulaziz \and *}
\institute{$^{1}$ Department of accounting, Faculty of Business Administration, Majmaah University, Al Majmaah 11952, Saudi Arabia}

\titlerunning{Journal of Statistics Applications \& Probability}
\authorrunning{An International Journal et al.}

%corresponding author email
\mail{j.abdelaziz@mu.edu.sa}

\received{----------. 2026}
\revised{----------. 2026}
\accepted{----------. 2026.}
\published{…. 2026.}

\abstracttext{This study aimed to identify the extent of commitment to Corporate Social Responsibility (CSR) accounting disclosure among family businesses in the Kingdom of Saudi Arabia. The research employed a descriptive-analytical approach, utilizing simple regression analysis. A questionnaire served as the primary instrument for data collection, distributed to a convenience sample of financial managers, accountants, and internal auditors within these companies, totaling (368) participants. Data were analyzed using various statistical methods via the SPSS software. The findings revealed a statistically significant level of commitment (at the significance level of α ≤ 0.05) by Saudi family businesses toward CSR accounting disclosure. In light of these results, the researchers proposed several recommendations, most notably: the necessity for the Saudi Organization for Chartered and Professional Accountants (SOCPA) to issue a dedicated standard or guiding framework for social responsibility, and the imperative for Saudi family businesses to integrate social information directly into their primary financial statements rather than limiting it to the supplementary notes.}

\keywords{Family Businesses, Accounting Disclosure, Kingdom of Saudi Arabia.} 



\maketitle


\par\vspace{12.0pt}\noindent {\large\bfseries 1. Introduction\par}\vspace{12.0pt}

The final decades of the twentieth century witnessed a growing emphasis on the social role of corporations, a shift largely driven by the retrenchment of government involvement in certain societal roles. Consequently, private sector companies have assumed a pivotal role in the developmental process. These entities have come to recognize that they do not operate in isolation from society, prompting a necessary expansion of their activities to address the needs of the communities in which they function. Today, corporations are viewed as economic entities operating within a social milieu; as social conditions improve, so too do the opportunities for these firms to thrive in a superior business environment \cite{ref1}, Currently, Corporate Social Responsibility (CSR) has gained paramount importance due to its profound impact on the surrounding community, serving as an effective instrument for driving progress and achieving social welfare.

Family businesses are among the most prevalent types of enterprises in emerging markets. \cite{ref8}, indicates that family firms represent two-thirds of businesses worldwide, contributing between 70\% and 90\% of the global annual GDP and providing 50\% to 80\% of jobs globally. Extensive research has focused on CSR within family businesses, which are characterized by family dynamics and a prevailing spirit of solidarity. These traits encourage cohesion among employees and the cultivation of long-term relationships with customers, suppliers, and other stakeholders \cite{ref6}, In the Kingdom of Saudi Arabia, family businesses constitute a cornerstone of various economic activities and a vital pillar of national production. They contribute between 60\% and 70\% of the national economy, with investments reaching SAR 250 billion, significantly influencing the GDP—a contribution the government remains keen to sustain \cite{ref7} .

\cite{ref2}, argue that CSR has imposed a new reality on corporations. Firms are no longer expected merely to utilize available resources for their own ends, but must also consider the requirements of the society of which they are an integral part. Accounting thought has responded to this reality in alignment with its mission to provide useful services to relevant stakeholders, leading to the emergence of "Social Responsibility Accounting."

Social Responsibility Accounting is a modern branch of the accounting field, holding significant local, regional, and international importance. It serves as a tool for accounting thought, providing essential information to parties interested in CSR by adopting an accounting model founded on social values \cite{ref3}, This model necessitates an expansion in accounting measurement and disclosure to encompass the external impacts of corporate activities, including social costs and social returns \cite{ref4}, Domestically, interest in Social Responsibility Accounting has surged following the rapid growth of the corporate sector across various sizes and legal forms. This necessitates monitoring corporate performance to ensure it meets societal needs at the highest levels. For these companies to ensure their continuity and growth, they must diligently attend to their social responsibilities toward various stakeholders, including shareholders, customers, employees, consumers, and the bbroader community \cite{ref5}, This study seeks to identify the extent of CSR disclosure among Saudi family businesses from the perspective of a sample of financial managers, accountants, and internal auditors within those firms.

\par\vspace{6.0pt}\noindent {\large\bfseries 2. Methodological Framework of the Study\par}\vspace{12.0pt}

\par\vspace{6.0pt}\noindent {\normalsize\bfseries 2.1. Study Problem\par}\vspace{6.0pt}

The perception of Corporate Social Responsibility (CSR) varies among family businesses; while some view it as an expenditure with no tangible returns, others utilize it as a strategic tool to enhance their corporate image within their operational environment. Commitment to implementing social responsibility accounting and disclosing it within financial reports and statements currently lacks legal enforceability, as it is often regarded as an ethical and moral obligation undertaken voluntarily by companies without government intervention. Numerous studies have addressed the challenges associated with CSR accounting disclosure across various legal corporate structures, concluding that there is a significant disparity in disclosure levels—both among companies within the same country and among those operating in the same sector across different nations \cite{ref37},.

From the two researchers' perspective, this variation is attributed to the absence of a dedicated accounting standard for social responsibility that provides certified and unified methods for measurement and disclosure. Furthermore, the application of social responsibility accounting remains essentially a voluntary and non-binding practice for corporations. Additionally, certain activities may have a detrimental impact on the environment and society, where disclosure could pose significant challenges for companies. Conversely, other firms whose activities do not harm the environment or society may overstate their positive social and environmental impacts \cite{ref35},. In light of this divergence in accounting disclosure levels, several knowledge gaps exist across various contexts that necessitate further investigation. Accordingly, the study problem can be articulated through the following primary question: To what extent do family businesses in the Kingdom of Saudi Arabia commit to CSR accounting disclosure?

\par\vspace{6.0pt}\noindent {\normalsize\bfseries 2.2. Significance of the Study\par}\vspace{6.0pt}

The theoretical significance of this study stems from the vital nature of its subject matter, as Corporate Social Responsibility (CSR) accounting disclosure has become a focal point for many corporations. This interest reflects a growing belief in their social role within their operational environments, alongside their traditional economic objectives. Furthermore, this study seeks to align with global efforts in the field of social responsibility accounting and contribute to highlighting the importance of disclosure in providing information that aids stakeholders in decision-making. This is particularly relevant given the scarcity of research within the Saudi and Arabic academic libraries regarding CSR accounting disclosure compared to studies conducted in Australia, the United States, and the United Kingdom; thus, this study represents an effort to enrich the literature on social responsibility accounting.

The practical significance is manifested in the two researchers' expectation that management within family businesses will benefit from the findings by raising awareness regarding the importance of CSR accounting disclosure. Additionally, the researchers hope that the results and recommendations will open new horizons for scholars and researchers to explore this topic through different perspectives and variables beyond those addressed in this study.

\par\vspace{6.0pt}\noindent {\normalsize\bfseries 2.3. Study Objectives\par}\vspace{6.0pt}

The study primarily aims to identify:

This study aims to identify the following:

\begin{enumerate}
  \item The extent of commitment among family businesses in the Kingdom of Saudi Arabia to accounting disclosure regarding environmental protection-related social responsibility.
  \item The extent of commitment among family businesses in the Kingdom of Saudi Arabia to accounting disclosure regarding community-related social responsibility.
  \item The extent of commitment among family businesses in the Kingdom of Saudi Arabia to accounting disclosure regarding employee-related social responsibility.
  \item d) The extent of commitment among family businesses in the Kingdom of Saudi Arabia to accounting disclosure regarding consumer protection-related social responsibility.
\end{enumerate}

\par\vspace{6.0pt}\noindent {\large\bfseries 3. Theoretical Framework of the Study\par}\vspace{12.0pt}

\par\medskip\noindent {\normalsize\bfseries 3.1. Family Businesses\par}\smallskip

Developing a unified and precise definition for the term "family business" remains challenging and lacks consensus in academic literature due to the overlap of various factors and the differing academic perspectives from which the term is viewed \cite{ref12}, The International Finance Corporation defined a family business as a company where the majority of votes are held by the controlling family \cite{ref38}, This concept emphasizes ownership and control as the primary distinguishing characteristics of a family firm. This definition aligns with the "Saudi Family Business Governance Guidance Note," which defines a family business as a company entirely owned or controlled by a specific family \cite{ref11}, \cite{ref13}, argue that a company is classified as a family business if it meets certain criteria, including that a controlling stake is owned by a single family, granting them the voting power to manage the company's operations and determine its future, whether directly or indirectly. Based on the foregoing, the researchers define a family business as the participation of more than one member of a single family in the ownership of the company's capital, while establishing controls that ensure the company's continuity in practice and control across generations.

The significance of family businesses lies in their substantial role in a nation's economic growth through their contribution to the GDP, as well as providing numerous employment opportunities and mitigating unemployment within societies.\cite{ref14}, \cite{ref15}. Family businesses are characterized by their continuity and intergenerational succession, possessing a distinct nature that sets them apart from other corporate entities. According to \cite{ref16}, family businesses are distinguished by a set of characteristics, including:

a) Legal Interlap between Ownership and Management: The phenomenon of overlapping ownership and management is one of the most prominent legal characteristics distinguishing family businesses. This differs from the principle of separation between ownership and management found in joint-stock companies, as family members occupy senior and middle management positions.

b) The Legal Nature of Family Values in the Corporate Organizational Structure: Family values serve as unwritten mandatory rules governing organizational behaviour in family firms. These values take various forms, ranging from ethical principles and family traditions to business practices inherited across generations.

C) Family Commitment: Family commitment and loyalty constitute a core element in shaping contractual relations within family businesses. Partner relations are built on mutual trust, imposing a moral and legal obligation of loyalty that transcends the rigid provisions of the company's articles of incorporation. Family businesses enjoy several advantages, which \cite{ref18}, \cite{ref17} have categorized as follows:

\begin{enumerate}
  \item Speed in decision-making and freedom from the formalities and procedures dictated by regulations and systems.
  \item The concentration of returns and profits within family members, which strengthens their financial position.
  \item Family businesses serve as incubators for creative ideas, providing funding, nurturing, and support; they also bolster social responsibility in the communities where they operate.
  \item They play a vital role in creating job opportunities within their environments, helping to solve or alleviate the problem of unemployment.
\end{enumerate}

The researchers believe that despite these advantages, family businesses are not without drawbacks and negatives. Most notably, the dominance of kinship over professional management, where the degree of relationship plays a role in the hiring process regardless of the candidate’s administrative or professional competence. Furthermore, there is often leniency regarding errors and a tendency to overlook some or all violations committed by family members working in the company. Additionally, there is a lack of clarity in defining authorities and responsibilities for family management members, stemming from a failure to apply sound standards for placing the right person in the right position, as well as the absence of formal job descriptions in this type of company.

\par\medskip\noindent {\large\bfseries 3.2. Corporate Social Responsibility (CSR)\par}\smallskip

Numerous definitions of Corporate Social Responsibility (CSR) have been proposed by private sector organizations, international bodies, authors, and researchers. However, its definition has undergone fundamental shifts over time, evolving alongside society and its expectations. Consequently, there is no universally accepted definition of CSR, nor is there a consensus on the specific social issues it encompasses \cite{ref5}.

CSR has been defined as a set of decisions and actions taken by an entity to achieve desired goals and prevailing societal values, representing a portion of the economic benefits the entity is expected to fulfill. In light of this, the researchers define CSR as the commitment of enterprises to achieve a positive impact on society and the environment by integrating ethical and social values into their business strategies \cite{ref6}; \cite{ref20}, ; \cite{ref40}, \cite{ref39}, Many researchers have sought to identify the dimensions of social responsibility, approaching them from various perspectives. The study by \cite{ref25}, indicates that CSR comprises four main dimensions:

a) The Social Dimension: Examples include donations, grants, and charitable social assistance that serve the community without a profit motive. Under this dimension, an entity can adopt a core community issue and work toward supporting and following up on its progress.

b) The Ethical Dimension: Also known as ethical responsibility, this dimension is based on moral principles and standards, as well as social norms and values. Within this framework, several indicators emerge, including equal employment opportunities, ethical considerations in consumption, adherence to human rights principles, and respect for the prevailing customs and traditions of society.

c) The Legal Dimension: This refers to legal responsibility and is based on principles of environmental protection, occupational safety, justice, and compliance with consumer laws. This dimension includes a wide range of elements intended for use by entities in a manner that enhances the relationship between the organization and stakeholders of diverse genders, races, and religions.

d) The Economic Dimension: This dimension implies that the entity must be economically viable and beneficial, striving to provide security for others by respecting fair competition rules and fully utilizing technological advancements to the extent that they do not cause harm to society or the environment.

\par\medskip\noindent {\normalsize\bfseries 3.3. Accounting Disclosure of Corporate Social Responsibility\par}\smallskip

\cite{ref36} , \cite{ref30}, define accounting disclosure as the process of providing financial and non-financial information related to an entity's financial and non-financial activities to various interested parties—such as investors, creditors, and regulatory bodies—through financial reports. \cite{ref9}, view accounting disclosure as the full and clear presentation of information contained within an entity's financial reports, provided that this information is characterized by integrity and transparency, and is delivered to users in a timely manner to facilitate rational decision-making. Based on the foregoing, the researchers define accounting disclosure as the method by which an entity informs society and its various stakeholders of its diverse activities through financial reports, in a manner that enables the measurement of its economic performance.

Accounting disclosure of Corporate Social Responsibility (CSR) is defined as the presentation of data and information related to an entity's social activities in a way that allows for the evaluation of its performance toward society \cite{ref21}, \cite{ref22}, defines it as the method or technique through which an entity informs the various segments of society about its activities with social implications via financial statements and reports. Accordingly, the researchers define CSR accounting disclosure as a methodology for measuring and communicating information resulting from an entity’s fulfilment of its social responsibilities, thereby enabling an assessment of its social performance.

The studies by \cite{ref23}, and Nebbe and Loualbia \cite{ref24}, sought to identify the significance of CSR accounting disclosure in corporations as follows:

\begin{enumerate}
  \item Determining the extent to which the entity fulfils and commits to its social responsibilities.
  \item Meeting society's needs for social information, which has become a fundamental requirement alongside financial information.
  \item Facilitating comparisons between different entities and across various time periods for the same entity.
  \item Developing disclosure standards by expanding the content and quality of disclosed information to include criteria and requirements that enable such disclosure.
\end{enumerate}

A number of studies, such as Tarshi and Yakhlif \cite{ref4},, indicate that three types of CSR accounting disclosure methods can be distinguished:

\begin{enumerate}
  \item Descriptive Reporting: This involves narrative explanations of the entity's social contributions and their resulting impacts. These reports are typically prepared when accountants face difficulties in quantifying activities.
  \item Social Cost Disclosure: This type of reporting discloses social costs but does not include the social benefits achieved, due to the difficulty of their measurement.
  \item Social Cost and Benefit Disclosure: These reports provide users with quantitative indicators of both the costs and benefits resulting from the social impact of the entity's activities.
\end{enumerate}

Based on the above, the researchers conclude that the methods and techniques of Corporate Social Responsibility (CSR) accounting disclosure vary. These include incorporating social performance within financial reports and statements or presenting them as separate reports and schedules. This variety clearly indicates the absence of a unified framework for presenting social responsibility information.

Furthermore, several factors and considerations may influence the level of CSR accounting disclosure in corporations. The study by Oneiza and Ali \cite{ref26}, identified these as follows:

\begin{enumerate}
  \item CSR disclosure entails additional costs for the entity resulting from the expansion of accounting disclosure practices.
  \item External auditors may oppose any expansion that imposes new responsibilities, particularly in areas requiring non-traditional auditing processes and activities.
  \item Media attention significantly impacts CSR accounting disclosure by highlighting undesirable corporate outputs. This encourages the community to demand that these companies mitigate their negative impacts on society.
\end{enumerate}

The researchers further contend that other factors also influence the level of an entity's social responsibility accounting disclosure, such as pressure from government agencies and the availability of expertise and qualified personnel capable of adopting such disclosure practices.

\par\vspace{6.0pt}\noindent {\large\bfseries 4. Previous Studies\par}\vspace{12.0pt}

The researchers have dedicated this section of the study to reviewing relevant previous literature, arranged chronologically from the most recent to the oldest, as follows:

\cite{ref29},: This study aimed to propose a model for reporting Corporate Social Responsibility (CSR) costs at the Libyan Iron and Steel Company. It also sought to evaluate the company's commitment to accounting disclosure of CSR activities and determine its overall social performance. Utilizing a descriptive-analytical approach, the study conducted personal interviews to analyse the level of measurement and accounting disclosure of social activity costs. The findings revealed that the company does not perform accounting measurement or disclosure of CSR costs within its financial statements and reports.

\cite{ref31},: This research aimed to highlight how industrial companies achieve competitive advantage through the dimensions of social responsibility, with a case study applied to Samsung from 2019 to 2022. Using a descriptive-analytical methodology, the study concluded that the company's mission and objectives align with societal values through its sponsorship of various cultural, social, and athletic activities.

\cite{ref28},: This theoretical study discussed the significance of CSR reporting and its future prospects amidst the global harmonization of financial reporting standards. The study evaluated three internationally recognized CSR reporting standards, reviewing the content of each to identify strengths, weaknesses, and their reliability as benchmarks for convergence and alignment. The results indicated that the Global Reporting Initiative (GRI) is the optimal standard for providing useful information for convergence without marginalizing other standards.

\cite{ref21},: This study aimed to identify the extent to which Islamic banks in Algeria implement social responsibility accounting and their level of disclosure. The research focused on Al Baraka Islamic Bank in Algeria using a descriptive-analytical approach. The findings indicated that Al Baraka Bank views social responsibility as a humanitarian aspect rather than a mandatory obligation; furthermore, the bank avoids disclosing information that might potentially harm stakeholders.

\cite{ref23},: This study aimed to assess the commitment of Libyan cement companies to disclosing CSR costs within their financial statements and reports. The researcher employed a descriptive-analytical approach, utilizing a questionnaire as the primary data collection instrument. The study concluded that the surveyed companies do not commit to disclosing CSR costs in their financial statements and reports.

\cite{ref10},: This study aimed to illustrate the methods of Corporate Social Responsibility (CSR) accounting disclosure within financial statements. The theoretical framework addressed the concept of social responsibility and its accounting treatments, while the empirical study focused on the SONATRACH corporation (Biskra branch) in Algeria. The research employed a content analysis methodology, examining the entity's financial statements for the year 2016. The findings indicated a deficiency in CSR-related accounting information within the financial statements, failing to meet user needs and resulting in a lack of data necessary to evaluate the social performance indicators of the corporation under study.

\cite{ref27}, : This research sought to identify the extent of CSR disclosure within the financial statements and reports of Saudi joint-stock companies. The study was applied to (170) companies, utilizing a questionnaire as the primary data collection instrument, which was distributed to a sample of accountants, auditors, and financial analysts. A key finding of the study was that Saudi joint-stock companies exhibit a high level of CSR disclosure in their financial reports concerning environmental protection, the local community, employees, and consumers.

\cite{ref5},: This study attempted to measure the disclosure of CSR elements by Saudi banks listed on the Saudi Stock Exchange (Tadawul). The research utilized an index comprising (16) elements distributed across various CSR domains. Employing a content analysis approach, the study analyzed the financial statements and reports of the target banks for the period 2012–2015. The findings revealed that the overall average CSR accounting disclosure index for Saudi banks was approximately 40\%. This percentage is considered below average, indicating that Saudi banks listed on the stock market do not disclose CSR elements to a high degree.

\cite{ref22}, : This study aimed to examine the experience of Danone Algeria in CSR accounting disclosure, as it is one of the leading companies in the production sector. The research adopted a descriptive methodology, using a questionnaire distributed to the company’s employees for data collection. The study concluded that the company discloses its social responsibility through financial statements, annual reports, sustainability reports, and socio-economic reports.

In reviewing the aforementioned literature, the researchers conclude that the findings of previous studies have been somewhat divergent. This study aligns with several previous works in its adoption of the descriptive-analytical approach and shares theoretical ground regarding family businesses and accounting disclosure. However, it distinguishes itself through its specific application environment—Saudi family businesses. During the literature review, the researchers did not encounter any study that addressed the current topic within this specific context, which underscores the originality and novelty of this research.

\par\vspace{6.0pt}\noindent {\large\bfseries 5. Study Hypotheses\par}\vspace{12.0pt}

In light of the study problem, its objectives, and the findings of previous literature, the study hypotheses can be formulated as follows:

\begin{enumerate}
  \item There is a statistically significant level of commitment (at the significance level of α ≤ 0.05) by family businesses in the Kingdom of Saudi Arabia toward accounting disclosure of social responsibility related to environmental protection.
  \item There is a statistically significant level of commitment (at the significance level of α ≤ 0.05) by family businesses in the Kingdom of Saudi Arabia toward accounting disclosure of social responsibility related to the local community.
  \item There is a statistically significant level of commitment (at the significance level of α ≤ 0.05) by family businesses in the Kingdom of Saudi Arabia toward accounting disclosure of social responsibility related to employees.
  \item There is a statistically significant level of commitment (at the significance level of α ≤ 0.05) by family businesses in the Kingdom of Saudi Arabia toward accounting disclosure of social responsibility related to consumer protection.
\end{enumerate}

\par\vspace{6.0pt}\noindent {\large\bfseries 6. Study Methodology\par}\vspace{12.0pt}

The current study adopted a descriptive-analytical approach, as it aligns with the nature of the research aimed at identifying the extent of Corporate Social Responsibility (CSR) disclosure among family businesses in the Kingdom of Saudi Arabia. This methodology is employed to describe the phenomenon under study and determine its characteristics, the relationships between its dimensions, and the factors influencing it, with the objective of reaching results that can be generalized to the study population \cite{ref32}.

\par\vspace{6.0pt}\noindent {\large\bfseries 7. Study Population and Sample\par}\vspace{12.0pt}

The study population consisted of all family businesses registered in the Kingdom of Saudi Arabia, totaling (278,680) companies (National Centre for Family Enterprises, 2025). According \cite{ref33}, a suitable sample size in descriptive research is 5\% of the study population if the population is very large. Based on this, the following sample size formula was applied:

\begin{equation}
	n = \frac{Z^{2} \cdot P \cdot (1 - P)}{d^{2}}
	= \frac{(1.96)^{2} \times 0.5 \times (1 - 0.5)}{(0.05)^{2}}
	= 384.16 \approx 384
\end{equation}

Where $n$ is the sample size, $Z$ is the standard normal score corresponding to the desired confidence level, $P$ is the estimated population proportion, and $d$ is the margin of error. At a 95\% confidence level ($Z = 1.96$), assuming $P = 0.5$ (to yield the maximum sample size) and a margin of error of $d = 0.05$, the required sample size ($n$) was calculated to be 384 companies. A total of 368 questionnaires were valid for statistical analysis, representing a response rate of 95.8\%. This response rate is considered high and acceptable for studies involving human participants, ensuring reliable results \citep{ref34}.
\par\vspace{6.0pt}\noindent {\large\bfseries 8. Study Instrument\par}\vspace{12.0pt}

The researchers developed a questionnaire to measure the extent of commitment among Saudi family businesses toward CSR accounting disclosure. After reviewing previous literature and research, the questionnaire was adapted to suit the current study's context. The instrument comprised two sections:

Section One: Included the personal data of the study sample, specifically: gender, educational qualification, years of experience, academic specialization, and job title.

Section Two: Consisted of four domains containing questions related to the study variables. This section included (41) items covering social responsibility dimensions and accounting disclosure. Responses were measured using a four-point Likert scale: (Strongly Agree (Very High) = 4, Agree (High) = 3, Disagree (Low) = 2, and Strongly Disagree (Very Low) = 1). To determine the interval length for each category, the range was calculated by subtracting the minimum value from the maximum value (4 - 1 = 3) and then dividing by the highest value in the scale (3 ÷ 4 = 0.74). This value was then added to the minimum value (1) to determine the upper limit of the category. The category intervals are detailed in Table (1) below.

Table (1): Four-point Likert Scale Classification (Mean Response Limits)

\JLTXaboveplacement \begin{jltxspan}%
\noindent\begin{minipage}{\linewidth}%
\centering
{\renewcommand{\arraystretch}{1.35}\small
\begin{tabular}{|>{\raggedright\arraybackslash}p{\dimexpr 0.3138\linewidth-2\tabcolsep-\arrayrulewidth\relax}|>{\centering\arraybackslash}p{\dimexpr 0.3417\linewidth-2\tabcolsep-\arrayrulewidth\relax}|>{\centering\arraybackslash}p{\dimexpr 0.2445\linewidth-2\tabcolsep-\arrayrulewidth\relax}|}
\hline
\textbf{Category} & \textbf{Lower Limit} & \textbf{Upper Limit} \\
\hline
Strongly Agree & 3.25 & 4.00 \\
\hline
Agree & 2.50 & 3.24 \\
\hline
Disagree & 1.75 & 2.49 \\
\hline
Strongly Disagree & 1.00 & 1.74 \\
\hline
\end{tabular}}
\captionof{table}{Table}%
\end{minipage}%
\end{jltxspan}\JLTXbelowplacement 

Prepared by the two researchers: Based on study data analysis, 2026.

\JLTXaboveplacement \begin{figure}[H]
\centering
\includegraphics[width=\linewidth,height=0.3\textheight,keepaspectratio]{media/office_obj_1.png}
\caption{Four-point Likert Scale Classification (Mean Response Limits)}
\label{fig:office1}
\end{figure}\JLTXbelowplacement 

The items in this section were distributed across four dimensions as follows:

First Dimension: Items related to CSR accounting disclosure regarding environmental protection (14 items).

Second Dimension: Items related to CSR accounting disclosure regarding the local community (10 items).

Third Dimension: Items related to CSR accounting disclosure regarding employees (8 items).

Fourth Dimension: Items related to CSR accounting disclosure regarding consumer protection (9 items).

Validity and Reliability of the Study Instrument

The validity of the study instrument (the questionnaire) refers to its ability to measure what it was intended to measure. The researchers assessed the validity and reliability of the instrument using two methods:

\begin{enumerate}
  \item Face Validity: The researchers presented the study instrument to a group of referees, including Professors and Associate Professors specialized in Accounting across several Saudi universities. Their expertise was sought to refine the tool, and their suggestions were incorporated to modify and finalize the questionnaire items.
  \item Cronbach’s Alpha Reliability: Following the necessary modifications and the subsequent distribution of the questionnaire to the study sample, the researchers collected the final data. Cronbach’s Alpha test was then employed to verify the reliability of the dimensions in the second section, as illustrated in Table (2).
\end{enumerate}

Table (2): Cronbach’s Alpha Coefficients for the Study Instrument

\JLTXaboveplacement \begin{jltxspan}%
\noindent\begin{minipage}{\linewidth}%
\centering
{\renewcommand{\arraystretch}{1.36}\small
\begin{tabular}{|>{\raggedright\arraybackslash}p{\dimexpr 0.3431\linewidth-2\tabcolsep-\arrayrulewidth\relax}|>{\centering\arraybackslash}p{\dimexpr 0.2520\linewidth-2\tabcolsep-\arrayrulewidth\relax}|>{\centering\arraybackslash}p{\dimexpr 0.4010\linewidth-2\tabcolsep-\arrayrulewidth\relax}|}
\hline
\textbf{Variable} & \textbf{Number of Items} & \textbf{Alpha Reliability Coefficient} \\
\hline
Environmental Protection & 14 & 0.71 \\
\hline
Local Community & 10 & 0.62 \\
\hline
Employees & 8 & 0.78 \\
\hline
Consumer Protection & 9 & 0.83 \\
\hline
\end{tabular}}
\captionof{table}{Table}%
\end{minipage}%
\end{jltxspan}\JLTXbelowplacement 

Prepared by the two researchers: Based on study data analysis, 2026.

\JLTXaboveplacement \begin{figure}[H]
\centering
\includegraphics[width=\linewidth,height=0.3\textheight,keepaspectratio]{media/office_obj_2.png}
\caption{Cronbach’s Alpha Coefficients for the Study Instrument}
\label{fig:office2}
\end{figure}\JLTXbelowplacement 

Table (2) demonstrates that the degree of consistency between the measurement items for each variable is robust, with values

ranging from 62\% to 83\%. Since the standard threshold for Cronbach’s Alpha reliability in administrative research is considered acceptable if it exceeds 60\%, the reliability coefficients for all study variables are deemed high and acceptable.

Statistical Methods Used in Data Analysis

The researchers utilized the Statistical Package for the Social Sciences (SPSS) to analyse the study data. Specifically:

\begin{enumerate}
  \item Cronbach’s Alpha Coefficient: To measure the reliability of the scales used.
  \item Kolmogorov-Smirnov Test: To verify whether the data follows a normal distribution.
  \item Descriptive Statistics: Mean and Standard Deviation were used to describe the study variables.
  \item Simple Regression Analysis: To test the validity of the study hypotheses.
\end{enumerate}

Description of Study Variables

The description of study variables aims to analyse the responses of the study sample regarding the items measuring each dimension. To achieve this, descriptive statistical methods were employed, including arithmetic means, standard deviations, significance levels, and degrees of agreement, as illustrated in Table (3).

Table (3): Description of Study Variables

\JLTXaboveplacement \begin{jltxspan}%
\noindent\begin{minipage}{\linewidth}%
\centering
{\small
\begin{tabular}{|>{\raggedright\arraybackslash}p{\dimexpr 0.0681\linewidth-2\tabcolsep-\arrayrulewidth\relax}|>{\raggedright\arraybackslash}p{\dimexpr 0.3001\linewidth-2\tabcolsep-\arrayrulewidth\relax}|>{\centering\arraybackslash}p{\dimexpr 0.0929\linewidth-2\tabcolsep-\arrayrulewidth\relax}|>{\centering\arraybackslash}p{\dimexpr 0.2490\linewidth-2\tabcolsep-\arrayrulewidth\relax}|>{\centering\arraybackslash}p{\dimexpr 0.2672\linewidth-2\tabcolsep-\arrayrulewidth\relax}|}
\hline
\textbf{NO} & \textbf{Variables} & \textbf{Mean} & \textbf{Standard Deviation} & \textbf{Degree of Agreement} \\
\hline
1 & Environmental Protection & 3.52 & 0.564 & Very High \\
\hline
2 & Local Community & 3.36 & 0.627 & Very High \\
\hline
3 & Employees & 3.23 & 0.459 & High \\
\hline
4 & Consumer Protection & 3.61 & 0.547 & Very High \\
\hline
5 & Accounting Disclosure & 3.46 & 0.621 & Very High \\
\hline
\multicolumn{2}{|l|}{\textbf{Total}} & \textbf{3.42} & \textbf{0.529} & \textbf{Very High} \\
\hline
\end{tabular}}
\captionof{table}{Table}%
\end{minipage}%
\end{jltxspan}\JLTXbelowplacement 

Prepared by the two researchers: Based on study data analysis, 2026.

\JLTXaboveplacement \begin{figure}[H]
\centering
\includegraphics[width=\linewidth,height=0.3\textheight,keepaspectratio]{media/office_obj_3.png}
\caption{Description of Study Variables}
\label{fig:office3}
\end{figure}\JLTXbelowplacement 

Table (3) demonstrates that the participants agreed with the items pertaining to Corporate Social Responsibility dimensions and accounting disclosure. The levels of agreement were relatively close, reaching a "Very High" degree for the variables (Environmental Protection, Local Community, Consumer Protection, and Accounting Disclosure), and a "High" degree for the (Employees) variable.

Normality Test

The researchers employed the Kolmogorov-Smirnov (K-S) test to verify the normality of the study data distribution, as shown in Table (4).

Table (4): Normality Test (One-Sample K-S)

\JLTXaboveplacement \begin{jltxspan}%
\noindent\begin{minipage}{\linewidth}%
\centering
{\small
\begin{tabular}{|>{\raggedright\arraybackslash}p{\dimexpr 0.2840\linewidth-2\tabcolsep-\arrayrulewidth\relax}|>{\centering\arraybackslash}p{\dimexpr 0.2085\linewidth-2\tabcolsep-\arrayrulewidth\relax}|>{\centering\arraybackslash}p{\dimexpr 0.1143\linewidth-2\tabcolsep-\arrayrulewidth\relax}|>{\centering\arraybackslash}p{\dimexpr 0.3892\linewidth-2\tabcolsep-\arrayrulewidth\relax}|}
\hline
\textbf{Variables} & \textbf{Number of Items} & \textbf{Z-Value} & \textbf{Significance Level (p-value)} \\
\hline
Environmental Protection & 14 & 0.642 & 1.435 \\
\hline
Local Community & 10 & 0.547 & 1.634 \\
\hline
Employees & 8 & 0.594 & 1.638 \\
\hline
Consumer Protection & 9 & 0.712 & 1.491 \\
\hline
\textbf{All Variables} & \textbf{41} & \textbf{0.645} & \textbf{1.523} \\
\hline
\end{tabular}}
\captionof{table}{Table}%
\end{minipage}%
\end{jltxspan}\JLTXbelowplacement 

Prepared by the two researchers: Based on study data analysis, 2026.

\JLTXaboveplacement \begin{figure}[H]
\centering
\includegraphics[width=\linewidth,height=0.3\textheight,keepaspectratio]{media/office_obj_4.png}
\caption{Normality Test (One-Sample K-S)}
\label{fig:office4}
\end{figure}\JLTXbelowplacement 

As indicated in Table (4), the calculated Z-values are lower than the critical Z-value (1.19), and the significance level (p-value) is greater than 0.05. This indicates that there are no statistically significant differences between the distribution of the questionnaire data and the normal distribution, confirming that the data follows a normal distribution.

Data Analysis and Discussion of Results

To test the validity of the study hypotheses, the researchers utilized simple linear regression analysis. The hypotheses proposed a statistically significant level of commitment (at the significance level of α ≤ 0.05) by family businesses in the Kingdom of Saudi Arabia toward CSR disclosure across four dimensions: Environmental Protection (X1), Local Community (X2), Employees (X3), and Consumer Protection (X4). The following sections detail the results of the data analysis and hypothesis testing.

Results of the First Hypothesis Testing

The first hypothesis states: "There is a statistically significant level of commitment (at the significance level of α ≤ 0.05) by family businesses in the Kingdom of Saudi Arabia toward accounting disclosure of social responsibility related to environmental protection." To verify this hypothesis, a simple regression equation was utilized, as illustrated in Table (5).

Table (5): Model Summary and ANOVA for the First Hypothesis

\JLTXaboveplacement \begin{jltxspan}%
\noindent\begin{minipage}{\linewidth}%
\centering
{\renewcommand{\arraystretch}{1.62}\small
\begin{tabular}{|>{\centering\arraybackslash}p{\dimexpr 0.1364\linewidth-2\tabcolsep-\arrayrulewidth\relax}|>{\centering\arraybackslash}p{\dimexpr 0.1347\linewidth-2\tabcolsep-\arrayrulewidth\relax}|>{\centering\arraybackslash}p{\dimexpr 0.1571\linewidth-2\tabcolsep-\arrayrulewidth\relax}|>{\centering\arraybackslash}p{\dimexpr 0.1000\linewidth-2\tabcolsep-\arrayrulewidth\relax}|>{\centering\arraybackslash}p{\dimexpr 0.1286\linewidth-2\tabcolsep-\arrayrulewidth\relax}|>{\centering\arraybackslash}p{\dimexpr 0.1202\linewidth-2\tabcolsep-\arrayrulewidth\relax}|>{\centering\arraybackslash}p{\dimexpr 0.1369\linewidth-2\tabcolsep-\arrayrulewidth\relax}|>{\centering\arraybackslash}p{\dimexpr 0.0674\linewidth-2\tabcolsep-\arrayrulewidth\relax}|}
\hline
\multirow{2}{*}{\textbf{Dependent Dimension}} & \multicolumn{4}{c|}{\textbf{Model Summary}} & \multicolumn{3}{c|}{\textbf{ANOVA}} \\
\cline{2-8}
 & \textbf{Correlation Coefficient (R)} & \textbf{Coefficient of Determination (R²)} & \textbf{Adjusted R²} & \textbf{Standard Error of the Estimate} & \textbf{Degrees of Freedom (df)} & \textbf{Calculated F-Value} & \textbf{Sig*} \\
\hline
Accounting Disclosure & 0.234 & 0.081 & 0.077 & 0.507 & 1 & 12.504 & 0.000 \\
\hline
\end{tabular}}
\captionof{table}{Table}%
\end{minipage}%
\end{jltxspan}\JLTXbelowplacement 

*Significant at α ≤ 0.05 level. Prepared by the two researchers: Based on study data analysis, 2026.

Table (5) demonstrates a positive correlation between the social responsibility dimension related to environmental protection (as the independent variable) and the dependent variable (accounting disclosure). The correlation coefficient (R) reached (0.234), while the coefficient of determination (R²) was (0.081). These results indicate that the environmental protection-related social responsibility dimension explains (8.1\%) of the variance in the accounting disclosure variable. The adjusted R² was (0.077), and the negligible difference (0.004) between R² and the adjusted R² confirms the model’s robust predictive capacity for the dependent variable. Furthermore, the table reveals the model's overall significance, with a calculated F-value of (12.504) and a significance level (sig.) of (0.000), which is less than (0.05). This result confirms a statistically significant commitment by Saudi family businesses toward CSR accounting disclosure related to environmental protection at (α ≤ 0.05) and (df=1). Consequently, the hypothesis is accepted.

Comparing this result with previous literature, it aligns with the findings of Ben Malik and Bougal (2024), Abdel Halim and Al-Nseirat (2018), and Qassoum (2018). Conversely, it diverges from the results of Mohamed and Al-Mazoughi (2025), Laila and Fatima (2022), Abu Stala (2022), Saadi et al. (2020), and Habbash and Al-Maliki (2018).

Results of the Second Hypothesis Testing

The second hypothesis states: "There is a statistically significant level of commitment (at the significance level of α ≤ 0.05) by family businesses in the Kingdom of Saudi Arabia toward accounting disclosure of social responsibility related to the local community." To verify this hypothesis, a simple regression equation was utilized, as illustrated in Table (6).

Table (6): Model Summary and ANOVA for the Second Hypothesis

\JLTXaboveplacement \begin{jltxspan}%
\noindent\begin{minipage}{\linewidth}%
\centering
{\renewcommand{\arraystretch}{1.62}\small
\begin{tabular}{|>{\centering\arraybackslash}p{\dimexpr 0.1364\linewidth-2\tabcolsep-\arrayrulewidth\relax}|>{\centering\arraybackslash}p{\dimexpr 0.1347\linewidth-2\tabcolsep-\arrayrulewidth\relax}|>{\centering\arraybackslash}p{\dimexpr 0.1571\linewidth-2\tabcolsep-\arrayrulewidth\relax}|>{\centering\arraybackslash}p{\dimexpr 0.1000\linewidth-2\tabcolsep-\arrayrulewidth\relax}|>{\centering\arraybackslash}p{\dimexpr 0.1286\linewidth-2\tabcolsep-\arrayrulewidth\relax}|>{\centering\arraybackslash}p{\dimexpr 0.1202\linewidth-2\tabcolsep-\arrayrulewidth\relax}|>{\centering\arraybackslash}p{\dimexpr 0.1369\linewidth-2\tabcolsep-\arrayrulewidth\relax}|>{\centering\arraybackslash}p{\dimexpr 0.0674\linewidth-2\tabcolsep-\arrayrulewidth\relax}|}
\hline
\multirow{2}{*}{\textbf{Dependent Dimension}} & \multicolumn{4}{c|}{\textbf{Model Summary}} & \multicolumn{3}{c|}{\textbf{ANOVA}} \\
\cline{2-8}
 & \textbf{Correlation Coefficient (R)} & \textbf{Coefficient of Determination (R²)} & \textbf{Adjusted R²} & \textbf{Standard Error of the Estimate} & \textbf{Degrees of Freedom (df)} & \textbf{Calculated F-Value} & \textbf{Sig*} \\
\hline
Accounting Disclosure & 0.273 & 0.078 & 0.073 & 0.419 & 1 & 14.367 & 0.000 \\
\hline
\end{tabular}}
\captionof{table}{Table}%
\end{minipage}%
\end{jltxspan}\JLTXbelowplacement 

*Significant at α ≤ 0.05 level. Prepared by the two researchers: Based on study data analysis, 2026.

Table (6) demonstrates a positive correlation between the social responsibility dimension related to the local community (as the independent variable) and the dependent variable (accounting disclosure). The correlation coefficient (R) reached (0.273), while the coefficient of determination (R²) was (0.078). These results indicate that the local community-related social responsibility dimension explains (7.8\%) of the variance in the accounting disclosure variable. The adjusted R² was (0.073), and the negligible difference (0.005) between R² and the adjusted R² confirms the model’s predictive capacity for the dependent variable. Furthermore, the table reveals the model's overall significance, with a calculated F-value of (14.367) and a significance level (sig.) of (0.000), which is less than (0.05). This result confirms a statistically significant commitment by Saudi family businesses toward CSR accounting disclosure related to the local community at (α ≤ 0.05) and (df=1). Consequently, the hypothesis is accepted.

Comparing this result with previous literature, it aligns with the findings of Ben Malik and Bougal (2024), Abdel Halim and Al-Nseirat (2018), and Qassoum (2018). Conversely, it diverges from the results of Mohamed and Al-Mazoughi (2025), Laila and Fatima (2022), Abu Stala (2022), Saadi et al. (2020), and Habbash and Al-Maliki (2018).

Results of the Third Hypothesis Testing

The third hypothesis states: "There is a statistically significant level of commitment (at the significance level of α ≤ 0.05) by family businesses in the Kingdom of Saudi Arabia toward accounting disclosure of social responsibility related to employees." To verify this hypothesis, a simple regression equation was utilized, as illustrated in Table (7).

Table (7): Model Summary and ANOVA for the Third Hypothesis

\JLTXaboveplacement \begin{jltxspan}%
\noindent\begin{minipage}{\linewidth}%
\centering
{\renewcommand{\arraystretch}{1.62}\small
\begin{tabular}{|>{\centering\arraybackslash}p{\dimexpr 0.1364\linewidth-2\tabcolsep-\arrayrulewidth\relax}|>{\centering\arraybackslash}p{\dimexpr 0.1347\linewidth-2\tabcolsep-\arrayrulewidth\relax}|>{\centering\arraybackslash}p{\dimexpr 0.1571\linewidth-2\tabcolsep-\arrayrulewidth\relax}|>{\centering\arraybackslash}p{\dimexpr 0.1000\linewidth-2\tabcolsep-\arrayrulewidth\relax}|>{\centering\arraybackslash}p{\dimexpr 0.1286\linewidth-2\tabcolsep-\arrayrulewidth\relax}|>{\centering\arraybackslash}p{\dimexpr 0.1202\linewidth-2\tabcolsep-\arrayrulewidth\relax}|>{\centering\arraybackslash}p{\dimexpr 0.1369\linewidth-2\tabcolsep-\arrayrulewidth\relax}|>{\centering\arraybackslash}p{\dimexpr 0.0674\linewidth-2\tabcolsep-\arrayrulewidth\relax}|}
\hline
\multirow{2}{*}{\textbf{Dependent Dimension}} & \multicolumn{4}{c|}{\textbf{Model Summary}} & \multicolumn{3}{c|}{\textbf{ANOVA}} \\
\cline{2-8}
 & \textbf{Correlation Coefficient (R)} & \textbf{Coefficient of Determination (R²)} & \textbf{Adjusted R²} & \textbf{Standard Error of the Estimate} & \textbf{Degrees of Freedom (df)} & \textbf{Calculated F-Value} & \textbf{Sig*} \\
\hline
Accounting Disclosure & 0.282 & 0.082 & 0.089 & 0.576 & 1 & 15.479 & 0.000 \\
\hline
\end{tabular}}
\captionof{table}{Table}%
\end{minipage}%
\end{jltxspan}\JLTXbelowplacement 

*Significant at α ≤ 0.05 level. Prepared by the two researchers: Based on study data analysis, 2026.

Table (7) demonstrates a positive correlation between the social responsibility dimension related to employees (as the independent variable) and the dependent variable (accounting disclosure). The correlation coefficient (R) reached (0.282), while the coefficient of determination (R²) was (0.082). These results indicate that the employee-related social responsibility dimension explains (8.2\%) of the variance in the accounting disclosure variable. The adjusted R² was (0.089), and the negligible difference (0.007) between R² and the adjusted R² confirms the model’s predictive capacity for the dependent variable. Furthermore, the table reveals the model's overall significance, with a calculated F-value of (15.479) and a significance level (sig.) of (0.000), which is less than (0.05). This result confirms a statistically significant commitment by Saudi family businesses toward CSR accounting disclosure related to employees at (α ≤ 0.05) and (df=1). Consequently, the hypothesis is accepted.

Comparing this result with previous literature, it aligns with the findings of Ben Malik and Bougal (2024), Abdel Halim and Al-Nseirat (2018), and Qassoum (2018). Conversely, it diverges from the results of Mohamed and Al-Mazoughi (2025), Laila and Fatima (2022), Abu Stala (2022), Saadi et al. (2020), and Habbash and Al-Maliki (2018).

Results of the Fourth Hypothesis Testing

The fourth hypothesis states: "There is a statistically significant level of commitment (at the significance level of α ≤ 0.05) by family businesses in the Kingdom of Saudi Arabia toward accounting disclosure of social responsibility related to consumer protection." To verify this hypothesis, a simple regression equation was utilized, as illustrated in Table (8).

Table (8): Model Summary and ANOVA for the Fourth Hypothesis

\JLTXaboveplacement \begin{jltxspan}%
\noindent\begin{minipage}{\linewidth}%
\centering
{\renewcommand{\arraystretch}{1.62}\small
\begin{tabular}{|>{\centering\arraybackslash}p{\dimexpr 0.1364\linewidth-2\tabcolsep-\arrayrulewidth\relax}|>{\centering\arraybackslash}p{\dimexpr 0.1347\linewidth-2\tabcolsep-\arrayrulewidth\relax}|>{\centering\arraybackslash}p{\dimexpr 0.1571\linewidth-2\tabcolsep-\arrayrulewidth\relax}|>{\centering\arraybackslash}p{\dimexpr 0.1000\linewidth-2\tabcolsep-\arrayrulewidth\relax}|>{\centering\arraybackslash}p{\dimexpr 0.1286\linewidth-2\tabcolsep-\arrayrulewidth\relax}|>{\centering\arraybackslash}p{\dimexpr 0.1202\linewidth-2\tabcolsep-\arrayrulewidth\relax}|>{\centering\arraybackslash}p{\dimexpr 0.1369\linewidth-2\tabcolsep-\arrayrulewidth\relax}|>{\centering\arraybackslash}p{\dimexpr 0.0674\linewidth-2\tabcolsep-\arrayrulewidth\relax}|}
\hline
\multirow{2}{*}{\textbf{Dependent Dimension}} & \multicolumn{4}{c|}{\textbf{Model Summary}} & \multicolumn{3}{c|}{\textbf{ANOVA}} \\
\cline{2-8}
 & \textbf{Correlation Coefficient (R)} & \textbf{Coefficient of Determination (R²)} & \textbf{Adjusted R²} & \textbf{Standard Error of the Estimate} & \textbf{Degrees of Freedom (df)} & \textbf{Calculated F-Value} & \textbf{Sig*} \\
\hline
Accounting Disclosure & 0.287 & 0.084 & 0.086 & 0.483 & 1 & 12.618 & 0.000 \\
\hline
\end{tabular}}
\captionof{table}{Table}%
\end{minipage}%
\end{jltxspan}\JLTXbelowplacement 

*Significant at α ≤ 0.05 level. Prepared by two the researchers: Based on study data analysis, 2026.

Table (8) demonstrates a positive correlation between the social responsibility dimension related to consumer protection (as the independent variable) and the dependent variable (accounting disclosure). The correlation coefficient (R) reached (0.287), while the coefficient of determination (R²) was (0.084). These results indicate that the consumer protection-related social responsibility dimension explains (8.4\%) of the variance in the accounting disclosure variable. The adjusted R² was (0.086), and the negligible difference (0.002) between R² and the adjusted R² confirms the model’s predictive capacity for the dependent variable. Furthermore, the table reveals the model's overall significance, with a calculated F-value of (12.618) and a significance level (sig.) of (0.000), which is less than (0.05). This result confirms a statistically significant commitment by Saudi family businesses toward CSR accounting disclosure related to consumer protection at (α ≤ 0.05) and (df=1). Consequently, the hypothesis is accepted.

Comparing this result with previous literature, it aligns with the findings of Ben Malik and Bougal (2024), Abdel Halim and Al-Nseirat (2018), and Qassoum (2018). Conversely, it diverges from the results of Mohamed and Al-Mazoughi (2025), Laila and Fatima (2022), Abu Stala (2022), Saadi et al. (2020), and Habbash and Al-Maliki (2018).

Results

\begin{enumerate}
  \item There is a statistically significant commitment (α ≤ 0.05) by family businesses in the Kingdom of Saudi Arabia toward accounting disclosure of social responsibility related to environmental protection. This is evidenced by disclosing costs related to soil pollution prevention, landscaping the surrounding environment, and costs associated with practices used to reduce environmental pollution.
  \item There is a statistically significant commitment (α ≤ 0.05) by family businesses in the Kingdom of Saudi Arabia toward accounting disclosure of social responsibility related to the local community. This is manifested through disclosing community donations for religious and social events, as well as disclosing costs for any obligations arising from legislative or contractual social requirements.
  \item There is a statistically significant commitment (α ≤ 0.05) by family businesses in the Kingdom of Saudi Arabia toward accounting disclosure of social responsibility related to employees. This is reflected in disclosing expenditures for creating a suitable work environment and disclosing support for employee subscriptions to various social activities.
  \item There is a statistically significant commitment (α ≤ 0.05) by family businesses in the Kingdom of Saudi Arabia toward accounting disclosure of social responsibility related to consumer protection. This is demonstrated through disclosing product development costs to meet customer needs and costs incurred by the companies to ensure compliance with quality standards.
\end{enumerate}

Recommendations

In light of the study’s findings, the researchers recommend the following:

\begin{enumerate}
  \item The Saudi Organization for Chartered and Professional Accountants (SOCPA) should issue a specialized standard or guiding manual for Corporate Social Responsibility. This should outline how family businesses should disclose CSR activities and unify the accounting disclosure methodology.
  \item Saudi family businesses should integrate social information directly into the primary financial statements, rather than limiting such data to the supplementary notes and disclosures.
  \item It is necessary to provide incentives for family businesses that demonstrate commitment to accounting disclosure of CSR across its various dimensions.
  \item Organizing training courses, conferences, and workshops focused on CSR for family businesses in Saudi Arabia to raise awareness regarding social responsibility accounting.
  \item Updating accounting curricula in Saudi higher education institutions by incorporating topics on social responsibility accounting and its disclosure methodologies.
\end{enumerate}

\begin{thebibliography}{99}
\bibitem{ref1}
Abdel Halim, A. H. M., \& Al-Nseirat, A. H. (2018). "The Extent of CSR Disclosure in the Financial Reports of Saudi Joint-Stock Companies: A Field Study." Journal of Social Sciences and Humanities, 51, 13–63.

\bibitem{ref2}
Abdel Salam, A., \& Bin Turkia, N. (2018). "Mechanisms of Accounting Measurement for Social Responsibility." International Journal of Islamic Economics, 69, 1–11.

\bibitem{ref3}
Abdel Haq, S. (2022). "Using Social Responsibility Accounting to Improve Employee Performance: A Case Study of the Social Affairs Directorate at SONATRACH – Biskra Agency." Unpublished PhD Thesis, University of Mohamed Khider, Biskra, Algeria.

\bibitem{ref4}
Tarshi, M., \& Yakhlif, I. (2017). "Disclosure of Social Responsibility from an Accounting Perspective." Al-Riyadah Journal for Business Economics, 3(2), 99–117.

\bibitem{ref5}
Habbash, M. S. M., \& Al-Maliki, A. M. (2018). "Accounting Disclosure of Social Responsibility in Saudi Banks." King Khalid University Journal for Humanities, 5(1), 175–216.

\bibitem{ref6}
Ali, A., \& Makiya, N. (2023). "Corporate Social Responsibility and Disclosure in Family Businesses: A Case Study of Al-Sindian Food Industries." Tishreen University Journal for Economic and Legal Sciences, 45(6), 1–24.

\bibitem{ref7}
Al-Muqaddam, M. S., \& Al-Azmi, A. O. M. (2024). "The Impact of Implementing Governance Dimensions on the Performance of Family Businesses in the Kingdom of Saudi Arabia." Journal of the Higher Institute for Qualitative Studies, 4(4), 1–53.

\bibitem{ref8}
Qarsham, A. A. (2023). "The Legal Form of Family Businesses." Journal of Legal and Economic Research, 85, 1–101.

\bibitem{ref9}
Malak, M. R., Ferial, A. M., \& Rashid, H. (2023). "Accounting Disclosure between the Financial Accounting System and International Accounting Standards." Paper presented at the International Scientific Conference: Requirements for Developing the Accounting System to Achieve Quality in Accounting Disclosure—Proposed Frameworks and Promising Prospects. Arab Democratic Center, Berlin, Germany, March 19, 2023.

\bibitem{ref10}
Saadi, A., Naima, Z., \& Aloun, M. L. (2020). "Accounting Disclosure of Social Responsibility in the Financial Statements of Economic Enterprises." Ara’ Journal for Economic and Administrative Studies, 2(2), 23–40.

\bibitem{ref11}
Saudi Family Business Governance Guidance Note and Model Charter." (2014). Ministry of Commerce and Industry, Kingdom of Saudi Arabia.

\bibitem{ref12}
Al-Zahrani, Y., \& Ahmed Al-Qasim. (2023). "The Legal Framework for Corporate Governance in the Kingdom of Saudi Arabia." Journal of Humanities and Administrative Sciences, 11(2), 1–39.

\bibitem{ref13}
Al-Owais, O. M. S., \& Abdel Hay, E. (2020). "Governance of Family Businesses in the United Arab Emirates: Challenges and Prospects: An Applied Legal Study." Journal of Sharia and Law, 82(34), 143–208.

\bibitem{ref14}
Fatiha, Q., Thabit, A., \& Abdel Latif, K. (2020). "Family Businesses: Models and Experiences." Collective Book, Abou Bakr Belkaïd University, Tlemcen, Algeria.

\bibitem{ref15}
Hamdi, M., Fallaq, S., \& Bouska, A. (2019). "Governance of Family Businesses: Mechanisms and Applications in the Nuqul Group - Jordan." Journal of Governance, Social Responsibility, and Sustainable Development, 1(1), 46–69.

\bibitem{ref16}
Al-Siyabi, H. K. M., \& Al-Hammadi, (2025) M. H. A. "Governance of Family Businesses in the Sultanate of Oman." Unpublished Master's Thesis, University of Ibra, Sultanate of Oman.

\bibitem{ref17}
Salman, S. T. (2011). "Strategic Planning for Family Businesses." Giza, Dar Al-Farouq for Cultural Investments, First Edition, Egypt.

\bibitem{ref18}
Darwish, N. A., \& Khanafseh, H. (2025). "The Reality of Succession Planning in Family Businesses in Palestine: A Case Study of Jerusalem and Ramallah Governorates." Arab Journal of Administrative Sciences, 32(1), 17–54.

\bibitem{ref19}
Hamed, K. A. (2014). "Dimensions of Social Responsibility to Enhance Creative Capabilities: An Applied Study on Selected Colleges at the University of Kufa." Journal of Economic and Administrative Sciences, 7(30), 1–21.

\bibitem{ref20}
Ismail, S. R. (2022). "The Importance of CSR Accounting Disclosure and its Impact on Achieving Sustainable Development: An Applied Study on the Saudi Paper Manufacturing Company in the Kingdom of Saudi Arabia." Journal of Economic, Administrative, and Legal Sciences, 7(12), 1–21.

\bibitem{ref21}
Laila, A., \& Fatima, S. (2022). "Accounting Disclosure of Social Responsibility in Islamic Banks: The Experience of Al Baraka Bank Algeria." Journal of Financial, Accounting, and Managerial Studies, 9(1), 1–20.

\bibitem{ref22}
Qassoum, H. (2018). "Disclosure of Social Responsibility in Economic Corporations: A Case Study of Danone." Al-Nadwa Journal for Legal Studies, 1–24.

\bibitem{ref23}
Abu Stala, A. I. (2022). "Accounting Disclosure of Social Responsibility in Financial Statements and Reports: An Applied Study on Libyan Cement Companies." Al-Asil Journal for Economic and Administrative Research, 6(1), 1–23.

\bibitem{ref24}
Nabou, M., \& Loualbia, F. (2021). "The Problem of Measurement and Disclosure of Social Responsibility in Economic Enterprises." Concepts Journal for Deep Philosophical and Human Studies, 9, 245–258.

\bibitem{ref25}
Saleh, M. A., Thajil, A. M., \& Mohamed, M. I. (2021). "Accounting Disclosure of Social Responsibility and its Impact on Achieving Sustainable Development Goals: An Applied Research at the General Company for Food Industries – Al-Diwaniyah Dairy Factory." Journal of Accounting and Financial Studies, 16(55), 1–20.

\bibitem{ref26}
Oniza, H. H., \& Ali, M. N. (2013). "The Impact of CSR Disclosure in Financial Statements on Users' Decisions: An Applied and Exploratory Study." Al-Ghary Journal of Economic and Administrative Sciences, 26, 153–189.

\bibitem{ref27}
Abdel Halim, A. M., \& Al-Nseirat, A. H. (2018). "The Extent of CSR Disclosure in the Financial Reports of Saudi Joint-Stock Companies: A Field Study." Journal of Humanities and Social Sciences, 51, 1–56.

\bibitem{ref28}
Mohamed, M. (2023). "Harmonizing Corporate Social Responsibility Disclosure Standards." Journal of Economy and Environment, 6(1), 1–23.

\bibitem{ref29}
Ahmed, K. B. M., \& Al-Mazoughi, O. M. (2025). "The Level of Accounting Measurement and Disclosure of CSR Activity Costs in Financial Statements: A Field Study on the Libyan Iron and Steel Company." Rimah Journal for Research and Studies, 116, 1–46.

\bibitem{ref30}
Qatati, N. A., Zoubik, F. F., \& Al-Mashat, A. A. (2013). "Disclosure of Social Responsibility Information in the Annual Reports of Commercial Banks Operating in Libya." Journal of Economic and Political Sciences, 2, 192–227.

\bibitem{ref31}
Ben Malik, A., \& Bougal, N. (2024). "The Problem of Achieving Competitive Advantage in Industrial Companies within the Context of Commitment to Social Responsibility Dimensions: A Case Study of Samsung for the Period 2019–2022." Jadeed Al-Iqtisad Journal, 19(1), 1–30.

\bibitem{ref32}
Al-Qahtani, S., Al-Amri, A., Al-Madhab, M., \& Al-Omar, B. (2013). Research Methodology in Behavioral Sciences. Riyadh: Obeikan Publishing.

\bibitem{ref33}
Mustafa, O. (2003). Scientific Research Skills. Center for Multidisciplinary Research and Studies.

\bibitem{ref34}
Babbie, E. R. (2020). The Practice of Social Research, 15th edition. Cengage Learning.

\bibitem{ref35}
Bayoud, N. \& Salem M. (2012). Corporate Social Responsibility Disclosure and Organizational Performance: The Case of Libya, A Mixed Methods Study, Doctor of Philosophy Dissertation, School of Accounting, Economics and Finance Faculty of Business and Law, University Of Southern Queensland.

\bibitem{ref36}
Blankespoor, E., deHaan, E., \& Marinovic, I. (2020). Disclosure processing costs, investors’ information choice, and equity market outcomes: A review. Journal of Accounting and Economics, 70(2-3), 101344.

\bibitem{ref37}
Hussainey, K., Elsayed, M., \& Razik, M. A. (2011). Factors affecting corporate social responsibility disclosure in Egypt. Journal of corporate social ownership and control, 8(4),432-443.

\bibitem{ref38}
International Financial Corporation. (2011). Family Business Governance Handbook, Op. cit,.p.12.

\bibitem{ref39}
Nugroho, H. (2021). Company Characteristics, Disclosure of Social Responsibility, and Its Impact on Company Performance: An Empirical Study in Indonesia. Journal of Asian Finance, Economics and Business, 8, (6), 889-895.

\bibitem{ref40}
Nyore, O., S, E-Obodoekwe, C., N, Pricilia, O., A, Chidiebere, N. (2021) , Impact of Environmental, Social and Governance Dimensions of Corporate Social Responsibility on Firm Performance: Evidence from Nigeria. Humanities and Social Sciences Letters, 9 (2), 220-236.

\end{thebibliography}


\end{document}
